% =====================================================================
%  Rozdział 3 - Środowisko badawcze (wersja robocza)
% =====================================================================
\chapter{Środowisko badawcze}\label{ch:srodowisko}

Rozdział opisuje stanowisko, na którym przeprowadzono wszystkie badania praktyczne: wymagania postawione środowisku, topologię sieci laboratoryjnej, konfigurację maszyn i~usług, platformę sprzętową wykorzystaną w~pomiarach wydajnościowych oraz przyjętą metodykę pomiarów. Opisana konfiguracja jest wspólna dla rozdziałów \ref{ch:ataki}--\ref{ch:obrona}; w~kolejnych rozdziałach przywoływane są jedynie elementy istotne dla danego eksperymentu.

% ---------------------------------------------------------------------
\section{Założenia i wymagania}\label{sec:wymagania}

Środowisko zaprojektowano tak, aby umożliwiało bezpieczne odtwarzanie badanych scenariuszy, rejestrację całego ruchu sieciowego oraz porównywanie mechanizmów w~wariancie podatnym i~poprawionym. Ze względu na dostępne zasoby całość musiała zmieścić się na jednym fizycznym komputerze, który pełni jednocześnie funkcję stanowiska do pomiarów wydajności kryptoanalizy z~wykorzystaniem procesora graficznego (ang. \ang{Graphics Processing Unit}, GPU). Sformułowane wymagania oraz sposób ich realizacji zestawiono w~tabeli~\ref{tab:wymagania}.

\begin{table}[htbp]
\centering\small
\begin{tabularx}{\textwidth}{@{}>{\raggedright\arraybackslash}p{4.2cm}L@{}}
\toprule
\textbf{Wymaganie} & \textbf{Sposób realizacji} \\
\midrule
Izolacja ruchu badawczego & ruch badawczy przesyłany wyłącznie w~wydzielonej sieci wewnętrznej hipernadzorcy, bez routingu do sieci zewnętrznej \\
Rozdzielenie ról & osobne maszyny wirtualne dla klienta, serwera usług i czujnika systemu detekcji włamań \\
Warianty porównawcze & każdy badany mechanizm uruchomiony równolegle w~wersji podatnej i~poprawionej na tym samym serwerze \\
Obserwowalność & czujnik z~interfejsem w~trybie nasłuchu (ang. \ang{promiscuous mode}) i~zapisem ruchu do plików \kod{pcap} \\
Powtarzalność & migawki stanu maszyn, wielokrotne powtarzanie pomiarów \\
Niezależność pomiarów wydajności & pomiary GPU wykonywane bezpośrednio w~systemie hosta przy wyłączonych maszynach wirtualnych oraz braku obciążenia innymi zadaniami \\

\bottomrule
\end{tabularx}
\caption{Wymagania stawiane środowisku badawczemu i~sposób ich realizacji}
\label{tab:wymagania}
\end{table}

% ---------------------------------------------------------------------
\section{Topologia stanowiska laboratoryjnego}\label{sec:topologia}

Laboratorium sieciowe zbudowano w~środowisku Oracle VirtualBox w~wersji 7.2.8~\cite{virtualbox}, uruchomionym na komputerze z~systemem Windows~11~Pro. Topologię stanowiska przedstawiono na rysunku~\ref{rys:topologia}. Składa się ono z~czterech maszyn wirtualnych połączonych wewnętrzną siecią VirtualBox o~nazwie \kod{labnet} (adresacja 10.10.10.0/24). Sieć ta nie posiada bramy ani połączenia z~siecią zewnętrzną -- cały ruch generowany w~trakcie badań pozostaje w~jej obrębie.

\begin{figure}[htbp]
\centering
\resizebox{\textwidth}{!}{%
\begin{tikzpicture}[
  font=\footnotesize,
  vm/.style={draw, rounded corners=3pt, minimum width=3.6cm, minimum height=3.2cm,
             align=center, fill=white, anchor=north},
  szyna/.style={line width=1pt},
  >={Stealth[length=2mm]}]
  % maszyny wirtualne (jednakowa wysokość, wyrównane krawędzie)
  \node[vm] (client) at (0,1.6)   {\textbf{Klient} (\kod{client})\\10.10.10.20\\skrypty pomiarowe};
  \node[vm] (server) at (4.2,1.6) {\textbf{Serwer} (\kod{server})\\10.10.10.30\\Docker:\\aplikacja testowa (:5000)\\nginx 1.27 (:443)\\OpenSSL 1.0.2u (:4433)};
  \node[vm, line width=1.2pt] (sensor) at (8.4,1.6) {\textbf{Czujnik} (\kod{sensor})\\10.10.10.40\\Snort 3, tcpdump};
  \node[font=\scriptsize, anchor=south] at (8.4,1.65) {tryb nasłuchu};
  % szyny
  \draw[szyna] (-1.95,-2.05) -- (10.35,-2.05);
  \draw[szyna, dashed] (-1.95,-2.8) -- (10.35,-2.8);
  \foreach \m in {client,server,sensor}{
    \draw ($(\m.south)+(-0.35,0)$) -- ($(\m.south)+(-0.35,-0.45)$) node[circle,fill,inner sep=1.3pt]{};
    \draw[dashed] ($(\m.south)+(0.35,0)$) -- ($(\m.south)+(0.35,-1.2)$) node[circle,fill,inner sep=1.3pt]{};
  }
  % host
  \node[draw, rounded corners=3pt, minimum width=12.3cm, minimum height=0.9cm, align=center,
        fill=black!4] (host) at (4.2,-4.3)
    {\textbf{Host}: Windows 11 Pro, VirtualBox 7.2.8 \quad$\cdot$\quad stanowisko GPU: Hashcat 7.1.2, NVIDIA GeForce GTX 1660};
  \draw[dashed,<->] (9.6,-2.8) -- (9.6,-3.85) node[midway,left,font=\scriptsize,align=right]{SSH z~hosta: 127.0.0.1:2220/2230/2240};
  % legenda
  \draw[szyna] (-1.9,-5.2) -- (-1.2,-5.2) node[right,font=\scriptsize]{sieć wewnętrzna \kod{labnet} 10.10.10.0/24 (karta \kod{enp0s3}) -- ruch badawczy};
  \draw[szyna,dashed] (-1.9,-5.6) -- (-1.2,-5.6) node[right,font=\scriptsize]{NAT VirtualBox (karta \kod{enp0s8}) -- instalacja pakietów i~zarządzanie};
\end{tikzpicture}}
\caption{Topologia stanowiska laboratoryjnego}
\label{rys:topologia}
\end{figure}

Każda maszyna posiada dwie karty sieciowe. Pierwsza, \kod{enp0s3}, jest podłączona do sieci \kod{labnet} i~ma statycznie przypisany adres; to przez nią przechodzi cały ruch, który wywołujemy podczas badań. Druga, \kod{enp0s8}, korzysta z~translacji NAT wbudowanej w VirtualBox i~służy wyłącznie do instalacji pakietów oraz zarządzania maszynami przez protokół SSH. Dostęp z~hosta realizowany jest przez przekierowanie portów 127.0.0.1:2220, 2230 i~2240 odpowiednio do maszyn klienta, serwera i~czujnika. Na karcie \kod{labnet} czujnika włączono tryb nasłuchu zarówno w~ustawieniach VirtualBox (zezwolenie na cały ruch), jak i~w~systemie operacyjnym, dzięki czemu czujnik rejestruje ruch wymieniany między pozostałymi maszynami, mimo że sam aktywnie w nim nie uczestniczy. Ustawienia karty sieciowej czujnika w~programie VirtualBox przedstawiono na rysunku~\ref{rys:sensor-siec}. Parametry maszyn zestawiono w~tabeli~\ref{tab:maszyny}.

\begin{figure}[htbp]
\centering
\includegraphics[width=0.85\textwidth]{rysunki/r3-sensor-siec.png}
\caption{Karta sieciowa czujnika: sieć wewnętrzna \kod{labnet} i~tryb nasłuchu (\ang{Promiscuous Mode}) ustawiony na zezwolenie na cały ruch}
\label{rys:sensor-siec}
\end{figure}

\begin{table}[htbp]
\centering\small
\begin{tabularx}{\textwidth}{@{}l l l c L@{}}
\toprule
\textbf{Maszyna} & \textbf{Adres} & \textbf{System} & \textbf{RAM / vCPU} & \textbf{Rola w~badaniach} \\
\midrule
\kod{client} & 10.10.10.20 & Ubuntu Server 24.04.5  & 1~GB / 1 & klient usług, uruchamianie skryptów pomiarowych \\
\kod{server} & 10.10.10.30 & Ubuntu Server 24.04.5  & 2~GB / 2 & usługi w~wariancie podatnym i~poprawionym \\
\kod{sensor} & 10.10.10.40 & Ubuntu Server 24.04.5  & 2~GB / 2 & czujnik systemu detekcji włamań, rejestracja ruchu \\
\bottomrule
\end{tabularx}
\caption{Maszyny wirtualne stanowiska laboratoryjnego}
\label{tab:maszyny}
\end{table}

Maszyny klienta, serwera i~czujnika utworzono jako klony połączone (ang. \ang{linked clones}) jednej maszyny bazowej \kod{ubuntu-base} z~systemem Ubuntu Server 24.04~\cite{ubuntuserver}. Klon połączony przechowuje jedynie różnice względem dysku maszyny bazowej, dlatego bezpośrednio po utworzeniu zajmuje kilkadziesiąt megabajtów zamiast kilku gigabajtów. Punktem wyjścia dla klonów jest migawka \kod{czysta-baza} maszyny bazowej. Po zainstalowaniu oprogramowania każdej z~maszyn wykonano migawkę \kod{baseline}, do której można wrócić przed powtórzeniem eksperymentu (rysunek~\ref{rys:migawki}). Budowę stanowiska zautomatyzowano: skrypt PowerShell uruchamiany na hoście tworzy klony i~konfiguruje karty sieciowe, a~skrypty uruchamiane w~maszynach nadają nazwy hostów, adresy w~sieci \kod{labnet} oraz instalują oprogramowanie właściwe dla roli maszyny.

\begin{figure}[htbp]
\centering
\includegraphics[width=0.85\textwidth]{rysunki/r3-virtualbox-migawki.png}
\caption{Maszyny laboratorium w~programie VirtualBox wraz z~migawką \kod{baseline}}
\label{rys:migawki}
\end{figure}

% ---------------------------------------------------------------------
\section{Konfiguracja maszyn i usług}\label{sec:konfiguracja}

\subsection{Serwer usług}\label{sec:serwer}

Wszystkie badane usługi uruchomiono na maszynie \kod{server} jako kontenery zarządzane narzędziem Docker Compose~\cite{docker}. Takie rozwiązanie pozwala uruchomić na jednym serwerze usługi o~skrajnie różnych wymaganiach -- w~tym serwer zbudowany na przestarzałej bibliotece kryptograficznej -- bez wzajemnego wpływu na konfigurację systemu. Definicję usług przedstawiono na listingu~\ref{lst:compose}, ich przeznaczenie w~tabeli~\ref{tab:uslugi}, a~stan uruchomionych kontenerów na rysunku~\ref{rys:docker}.

\begin{figure}[htbp]
\centering
\includegraphics[width=\textwidth]{rysunki/r3-server-docker.png}
\caption{Kontenery usług uruchomione na maszynie \kod{server} (polecenie \kod{docker compose ps})}
\label{rys:docker}
\end{figure}

\begin{lstlisting}[float=htbp,language={},caption={Definicja usług serwera laboratoryjnego (plik \kod{docker-compose.yml}, fragment)},label={lst:compose}]
services:
  webapp:                       # aplikacja testowa (HTTP)
    build: ./webapp
    ports: ["5000:5000"]
  modern-tls:                   # serwer TLS - konfiguracja utwardzona
    image: nginx:1.27-alpine
    volumes:
      - ./nginx-modern/nginx.conf:/etc/nginx/nginx.conf:ro
      - ./certs:/certs:ro
    ports: ["443:443"]
  legacy-tls:                   # serwer TLS - konfiguracja przestarzała
    build: ./legacy-tls         # OpenSSL 1.0.2u kompilowany ze źródeł
    volumes: ["./certs:/certs:ro"]
    ports: ["4433:4433"]
\end{lstlisting}

\begin{table}[htbp]
\centering\small
\begin{tabularx}{\textwidth}{@{}l l L@{}}
\toprule
\textbf{Port} & \textbf{Usługa} & \textbf{Przeznaczenie} \\
\midrule
5000/tcp & aplikacja testowa (Flask) & ciasteczka sesyjne i~podpisy żądań w~wariancie podatnym i~poprawionym (rozdz.~\ref{ch:ataki}) \\
443/tcp  & nginx 1.27 & serwer TLS w~konfiguracji utwardzonej: tylko TLS~1.2 i~1.3 oraz zestawy szyfrów z~uwierzytelnieniem \\
4433/tcp & OpenSSL 1.0.2u (\kod{s\_server}) & serwer TLS w~konfiguracji przestarzałej: protokoły od SSLv3 do TLS~1.2, w~tym zestawy z~trybem CBC i~szyfrem RC4 \\
\bottomrule
\end{tabularx}
\caption{Usługi uruchomione na serwerze laboratoryjnym}
\label{tab:uslugi}
\end{table}

Aplikację testową napisano w~języku Python z~użyciem szkieletu Flask~\cite{flask} i~biblioteki kryptograficznej PyCryptodome~\cite{pycryptodome}. Aplikacja udostępnia trzy pary mechanizmów, w~których wariant podatny i~poprawiony różnią się wyłącznie zastosowaną konstrukcją kryptograficzną (tabela~\ref{tab:aplikacja}). Klucze szyfrujące (128-bitowe, osobne dla każdego wariantu) oraz hasło używane do podpisywania żądań generowane są losowo przy każdym uruchomieniu aplikacji. Aplikację udostępniono celowo bez TLS, aby jej ruch był w~całości widoczny dla sensora. Treść ciasteczka sesyjnego ma postać \kod{user=<nazwa>;admin=0;uid=1001}; pole \kod{admin} decyduje o~dostępie do funkcji administracyjnych, co pozwala ocenić nie tylko to, czy modyfikacja została przyjęta, ale też jaki ma skutek.

\begin{table}[htbp]
\centering\small
\begin{tabularx}{\textwidth}{@{}>{\raggedright\arraybackslash}p{2.9cm}LLc@{}}
\toprule
\textbf{Mechanizm} & \textbf{Wariant podatny} & \textbf{Wariant poprawiony} & \textbf{Podrozdz.} \\
\midrule
ciasteczko sesyjne (szyfr strumieniowy) & AES-128-CTR bez kodu uwierzytelniającego; \kod{/ctr/login}, \kod{/ctr/profile} & AES-128-GCM; \kod{/secure/login}, \kod{/secure/profile} & \ref{sec:ctr} \\
ciasteczko sesyjne (szyfr blokowy) & AES-128-CBC bez kodu uwierzytelniającego; \kod{/login}, \kod{/profile} & AES-128-GCM (jw.) & \ref{sec:cbc} \\
podpis żądania API & SHA-256(sekret~$\|$~wiadomość); \kod{/api/sample}, \kod{/api/transfer} & HMAC-SHA-256; \kod{/api/secure/sample}, \kod{/api/secure/transfer} & \ref{sec:lenext} \\
\bottomrule
\end{tabularx}
\caption{Pary mechanizmów aplikacji testowej}
\label{tab:aplikacja}
\end{table}

Do badania negocjacji parametrów sesji przygotowano dwa serwery TLS posługujące się tym samym, samopodpisanym certyfikatem (klucz RSA 2048~bitów, nazwa \kod{lab-server}). Serwer w~konfiguracji utwardzonej to nginx~\cite{nginxssl}, w~którym dopuszczono wyłącznie protokoły TLS~1.2 i~1.3 oraz zestawy szyfrów z~uwierzytelnieniem (listing~\ref{lst:nginx}). Serwer w~konfiguracji przestarzałej to program \kod{s\_server} z~biblioteki OpenSSL~\cite{openssl} w~wersji 1.0.2u. Bibliotekę tą uruchomiono z możliwie szeroką listą zestawów szyfrów (listing~\ref{lst:legacy}).

\begin{lstlisting}[float=htbp,language={},caption={Konfiguracja protokołów i~szyfrów serwera utwardzonego (fragment)},label={lst:nginx}]
listen 443 ssl;
ssl_certificate     /certs/server.crt;
ssl_certificate_key /certs/server.key;
ssl_protocols TLSv1.2 TLSv1.3;
ssl_ciphers ECDHE-ECDSA-AES128-GCM-SHA256:ECDHE-RSA-AES128-GCM-SHA256:
            ECDHE-RSA-AES256-GCM-SHA384:ECDHE-RSA-CHACHA20-POLY1305;
ssl_prefer_server_ciphers on;
\end{lstlisting}

\begin{lstlisting}[float=htbp,language={},caption={Uruchomienie serwera w~konfiguracji przestarzałej},label={lst:legacy}]
exec /opt/openssl-legacy/bin/openssl s_server \
  -accept 4433 \
  -cert /certs/server.crt -key /certs/server.key \
  -cipher 'ALL:RC4:!aNULL:!eNULL' \
  -www
\end{lstlisting}

\subsection{Maszyna klienta}\label{sec:klient}

Maszyna \kod{client} służy jako klient wszystkich badanych usług i~miejsce uruchamiania skryptów pomiarowych. Oprócz systemowej biblioteki OpenSSL w~wersji 3.0.13 zainstalowano na niej równolegle bibliotekę OpenSSL 1.0.2u. Konieczne było takowe rozwiązanie, ponieważ współczesne wersje biblioteki nie pozwalają już nawiązać połączenia z~użyciem protokołów SSLv3, TLS~1.0 i~TLS~1.1, co uniemożliwiłoby zbadanie, które z~tych wersji akceptuje serwer. Skrypty pomiarowe uruchamiane na kliencie napisano w~języku Python (wersja 3.12.3) z~wykorzystaniem wyłącznie biblioteki standardowej oraz w~języku powłoki Bash.

\subsection{Czujnik systemu detekcji włamań}\label{sec:czujnik}

Na maszynie \kod{sensor} zainstalowano IDS (ang. \ang{Intrusion Detection System}) Snort w~wersji 3.12.2.0~\cite{snort3}, skompilowany ze źródeł. Na czas kompilacji przydział pamięci maszyny zwiększono do 4~GB. Konfiguracja czujnika składa się z~pliku bazowego \kod{snort.lua}, pliku z~ustawieniami laboratorium \kod{lab.lua} (wczytywanego opcją \kod{--tweaks lab}) oraz pliku reguł własnych \kod{local.rules}. W~pliku \kod{lab.lua} włączono inspektor \kod{arp\_spoof} i~wpisano do niego poprawne powiązania adresów IP z~adresami sprzętowymi kart sieciowych klienta i~serwera, odczytane z~tablicy sąsiadów czujnika. Szczegóły konfiguracji detekcji opisano w~podrozdziale~\ref{sec:ids}. Do zapisu ruchu w~plikach \kod{pcap} wykorzystywano program tcpdump~\cite{tcpdump}. Zapisane pliki analizowano następnie offline, co pozwala wielokrotnie uruchamiać detekcję na identycznym materiale.


\subsection{Wersje oprogramowania}

Wersje oprogramowania wszystkich maszyn zestawiono w~tabeli~\ref{tab:wersje}. Wykaz sporządzono po zakończeniu instalacji i~nie zmieniano go w~trakcie badań.

\begin{table}[htbp]
\centering\small
\begin{tabularx}{\textwidth}{@{}l L l@{}}
\toprule
\textbf{Maszyna} & \textbf{Oprogramowanie} & \textbf{Wersja} \\
\midrule
\kod{server}, \kod{client}, \kod{sensor} & Ubuntu Server / jądro Linux & 24.04.5 LTS / 6.8.0-147 \\
\kod{server} & Docker / Docker Compose & 29.1.3 / 2.40.3 \\
\kod{server} & nginx (obraz kontenera) & 1.27 \\
\kod{server}, \kod{client} & OpenSSL (instalacja przestarzała) & 1.0.2u \\
\kod{client} & OpenSSL (systemowy) / Python & 3.0.13 / 3.12.3 \\
\kod{sensor} & Snort & 3.12.2.0 \\

host & Windows / VirtualBox & 11 Pro / 7.2.8 \\
\bottomrule
\end{tabularx}
\caption{Wersje oprogramowania stanowiska laboratoryjnego}
\label{tab:wersje}
\end{table}

% ---------------------------------------------------------------------
\section{Platforma sprzętowa do badań wydajnościowych}\label{sec:platforma}

Badania wydajności kryptoanalizy offline (rozdział~\ref{ch:gpu}) przeprowadzono na tym samym komputerze, na którym działa laboratorium sieciowe, ale bezpośrednio w~systemie Windows, a~nie w~maszynach wirtualnych. Parametry zestawu przedstawiono w~tabeli~\ref{tab:platforma}.

\begin{table}[htbp]
\centering\small
\begin{tabularx}{\textwidth}{@{}>{\raggedright\arraybackslash}p{4.3cm}L@{}}
\toprule
\textbf{Element} & \textbf{Parametry} \\
\midrule
Procesor (CPU) & AMD Ryzen 5 5600, 6 rdzeni, 12 wątków \\
Pamięć operacyjna & 16~GB \\
Karta graficzna (GPU) & NVIDIA GeForce GTX 1660, 6~GB pamięci GDDR5 \\
System operacyjny & Windows 11 Pro \\
Sterownik karty graficznej & NVIDIA 610.88 (CUDA 13.3) \\
Narzędzie do kryptoanalizy & Hashcat 7.1.2 \\
Obliczenia na GPU & OpenCL 3.0 (platforma NVIDIA CUDA w~sterowniku) \\
Obliczenia na CPU & OpenCL 3.0 (Intel CPU Runtime for OpenCL) \\
Pliki testowe & VeraCrypt 1.26.29, KeePassXC 2.8.0, 7-Zip \\
Generowanie skrótów Argon2 & Python 3.12, biblioteka \kod{argon2-cffi} \\
\bottomrule
\end{tabularx}
\caption{Platforma sprzętowa i~programowa badań wydajnościowych}
\label{tab:platforma}
\end{table}

Do pomiarów wykorzystano narzędzie Hashcat~\cite{hashcat}, które pozwala uruchomić te same algorytmy zarówno na karcie graficznej, jak i~na procesorze. Oba urządzenia obsługiwane były przez interfejs OpenCL: karta graficzna -- przez implementację dostarczaną ze sterownikiem NVIDIA, procesor -- przez środowisko Intel CPU Runtime for OpenCL. Natywny moduł CUDA Hashcata nie został uruchomiony, ponieważ wymaga bibliotek pakietu CUDA Toolkit, których nie instalowano; program automatycznie przełączył się wtedy na OpenCL. Listę urządzeń wykrytych przez Hashcat przedstawiono na rysunku~\ref{rys:hashcat-urz}. Istotny dla dalszych wyników jest wiersz \kod{Memory.Total} karty graficznej: z~6143~MB pamięci interfejs OpenCL pozwala przydzielić w~jednym bloku najwyżej 1535~MB.

\begin{figure}[htbp]
\centering
\includegraphics[width=0.62\textwidth]{rysunki/r3-hashcat-urzadzenia.png}
\caption{Urządzenia obliczeniowe wykryte przez Hashcat (polecenie \kod{hashcat -I})}
\label{rys:hashcat-urz}
\end{figure}

Pliki testowe, czyli kontenery szyfrowane, bazy haseł i~zaszyfrowane archiwa, utworzono w~programach VeraCrypt~\cite{veracrypt}, KeePassXC~\cite{keepassxc} oraz 7-Zip, a~dane potrzebne do ataku offline wyodrębniono z~nich pomocniczymi narzędziami \kod{7z2hashcat} i~\kod{keepass2john}.

Na czas pomiarów wydajności wszystkie maszyny wirtualne oraz wszystkie procesy mogące powodować zakrzywnienie wyników były wyłączone. W~trakcie pomiarów na karcie graficznej co 2~sekundy zapisywano jej temperaturę, taktowanie, pobór mocy i~obciążenie (program \kod{nvidia-smi}). Pozwoliło to wykluczyć wpływ parametrów urządzenia, np. temperatury na wyniki.

% ---------------------------------------------------------------------
\section{Metodyka pomiarów}\label{sec:metodyka}

Wszystkie eksperymenty przeprowadzono według wspólnego schematu, przedstawionego na rysunku~\ref{rys:przebieg}. Każdy pomiar poprzedzało przygotowanie środowiska: uruchomienie potrzebnych maszyn i~usług, a~w~przypadku badań ruchu sieciowego -- rozpoczęcie jego zapisu na maszynie czujnik. Sam pomiar wykonywał skrypt, który zapisywał zarówno wynik zbiorczy, jak i~surowe dane każdej próby. Wyniki weryfikowano, zapisywano i~opracowywano w~postaci tabel i~wykresów.

\begin{figure}[htbp]
\centering
\begin{tikzpicture}[
  font=\footnotesize,
  krok/.style={draw, rounded corners=3pt, minimum width=2.5cm, minimum height=1.0cm,
               align=center, fill=black!4},
  opis/.style={align=center, text width=2.6cm, font=\scriptsize},
  >={Stealth[length=2mm]}]
  \node[krok] (k1) at (0,0)     {\textbf{Przygotowanie}};
  \node[krok] (k2) at (3.05,0)  {\textbf{Pomiar}};
  \node[krok] (k3) at (6.1,0)   {\textbf{Rejestracja}};
  \node[krok] (k4) at (9.15,0)  {\textbf{Weryfikacja}};
  \node[krok] (k5) at (12.2,0)  {\textbf{Opracowanie}};
  \foreach \a/\b in {k1/k2,k2/k3,k3/k4,k4/k5} \draw[->] (\a) -- (\b);
  \node[opis, below=2mm of k1] {migawka \kod{baseline}, uruchomienie usług, zapis ruchu na czujniku};
  \node[opis, below=2mm of k2] {skrypt pomiarowy, stałe ziarno, seria prób};
  \node[opis, below=2mm of k3] {pliki CSV i~TXT, pliki \kod{pcap}, log karty graficznej};
  \node[opis, below=2mm of k4] {powtórzenia, porównanie z~wartością oczekiwaną};
  \node[opis, below=2mm of k5] {archiwum wyników, tabele i~wykresy};
\end{tikzpicture}
\caption{Schemat przebiegu pojedynczego eksperymentu}
\label{rys:przebieg}
\end{figure}

Przyjęto następujące założenia:
\begin{itemize}
  \item \textbf{Porównanie wariantów.} Każdy mechanizm badano w~wariancie podatnym i~poprawionym w~identycznych warunkach, tym samym skryptem. Dzięki temu różnica w~wynikach wynika wyłącznie z~zastosowanej konstrukcji lub konfiguracji.
  \item \textbf{Automatyzacja.} Pomiary wykonywały skrypty zestawione w~tabeli~\ref{tab:skrypty}. Skrypty zapisywały surowe dane każdej próby, co umożliwiało późniejsze sprawdzenie wyniku zbiorczego.
  \item \textbf{Powtarzalność.} Serie losowych prób wykonywano ze stałym ziarnem generatora pseudolosowego (20261025 dla rozdziału~\ref{ch:ataki}, 20261029 dla zbioru testowego w~rozdziale~\ref{ch:obrona}), dzięki czemu można je dokładnie powtórzyć. Pomiary wydajności powtarzano trzykrotnie (dla funkcji Argon2 z~pamięcią 64 i~256~MiB -- dwukrotnie) i~podawano wartość średnią.
  \item \textbf{Weryfikacja poprawności wyników.} Przed uruchomieniem na stanowisku każdy skrypt sprawdzano na lokalnej kopii aplikacji testowej. Tam, gdzie było to możliwe, wynik pomiaru porównywano z~wartością wynikającą z~budowy badanego mechanizmu.
  \item \textbf{Analiza offline.} Ruch sieciowy zapisywano do plików \kod{pcap}, a~detekcję uruchamiano na zapisanym ruchu. Pozwala to porównywać różne wersje reguł na identycznym materiale wejściowym.
  \item \textbf{Archiwizacja.} Wszystkie wyniki (pliki CSV i~TXT, pliki \kod{pcap}, logi) dostarczono do~repozytorium projektu wraz ze skryptami, które je wytworzyły.
\end{itemize}

\begin{table}[htbp]
\centering\small
\begin{tabularx}{\textwidth}{@{}>{\raggedright\arraybackslash}p{5.2cm}l L c@{}}
\toprule
\textbf{Skrypt} & \textbf{Maszyna} & \textbf{Przeznaczenie} & \textbf{Rozdz.} \\
\midrule
\kod{r4-skan-tls.sh}, \kod{r4-scsv.sh} & klient & wersje protokołu akceptowane przez serwery TLS, reakcja na sygnał obniżenia wersji & \ref{ch:ataki} \\
\kod{r4-pomiar-integralnosci.py} & klient & odporność ciasteczek i~podpisów na modyfikację & \ref{ch:ataki} \\
\kod{pomiar-bench.ps1} & host & test wydajności GPU z~zapisem parametrów karty & \ref{ch:gpu} \\
\kod{pomiar-veracrypt.ps1}, \kod{pomiar-hash.ps1} & host & prędkość ataku na własne pliki testowe & \ref{ch:gpu} \\
\kod{pomiar-gpu-cpu.ps1} & host & porównanie GPU i~CPU dla wybranych algorytmów & \ref{ch:gpu} \\
\kod{gen-argon2.py}, \kod{pomiar-argon2.ps1} & host & skalowanie funkcji Argon2 względem pamięci & \ref{ch:gpu} \\
\kod{local.rules}, \kod{lab.lua} & czujnik & reguły i~konfiguracja detekcji & \ref{ch:obrona} \\
\kod{r6-gen-zbior.py}, \kod{r6-entropia.py}, \kod{r6-realne.py} & host & zbiór testowy i~klasyfikacja plików zaszyfrowanych & \ref{ch:obrona} \\
\bottomrule
\end{tabularx}
\caption{Skrypty i~pliki konfiguracyjne wykorzystane w~pomiarach}
\label{tab:skrypty}
\end{table}
